\documentclass[11pt,a4paper]{ctexart}
%\usepackage{CJKutf8}
\usepackage{beamerthemesplit}
\usepackage{}
%\usetheme{Frankfurt}
%\usetheme{CambridgeUS}
%\usecolortheme{beaver}
%\usetheme{AnnArbor}
\usetheme{Dresden}
%\usebeamercolor{beetle}
\usepackage{xeCJK}
\setCJKmainfont{AR PL KaitiM GB}
%\useoutertheme{miniframes}
\usepackage{amsmath}
\usepackage{graphicx}
\usepackage{float} 
\usepackage{subfigure}
\usepackage{amssymb}
\usepackage{graphicx}
\usepackage{eufrak}
\usepackage{color}
\usepackage{array}
\usepackage{slashed}
\usepackage{simplewick}
\usepackage{tikz}
\usepackage{tcolorbox}
\usepackage[T1]{fontenc}
\graphicspath{{../figures/}}

%%figures
\def\lfig#1#2{\includegraphics[width=#1 in]{#2}}
\def\addfig#1#2{\begin{center}\includegraphics[width=#1 in]{#2}\end{center}}
\def\wulian{\includegraphics[width=0.18in]{emoji_wulian.jpg}}
\def\bigwulian{\includegraphics[width=0.35in]{emoji_wulian.jpg}}
\def\bye{\includegraphics[width=0.18in]{emoji_bye.jpg}}
\def\bigbye{\includegraphics[width=0.35in]{emoji_bye.jpg}}
\def\huaixiao{\includegraphics[width=0.18in]{emoji_huaixiao.jpg}}
\def\bighuaixiao{\includegraphics[width=0.35in]{emoji_huaixiao.jpg}}
\def\jianxiao{\includegraphics[width=0.18in]{emoji_jianxiao.jpg}}
\def\bigjianxiao{\includegraphics[width=0.35in]{emoji_jianxiao.jpg}}
%% colors
\def\blacktext#1{{\color{black}#1}}
\def\bluetext#1{{\color{blue}#1}}
\def\redtext#1{{\color{red}#1}}
\def\darkbluetext#1{{\color[rgb]{0,0.2,0.6}#1}}
\def\skybluetext#1{{\color[rgb]{0.2,0.7,1.}#1}}
\def\cyantext#1{{\color[rgb]{0.,0.5,0.5}#1}}
\def\greentext#1{{\color[rgb]{0,0.7,0.1}#1}}
\def\darkgray{\color[rgb]{0.2,0.2,0.2}}
\def\lightgray{\color[rgb]{0.6,0.6,0.6}}
\def\gray{\color[rgb]{0.4,0.4,0.4}}
\def\blue{\color{blue}}
\def\red{\color{red}}
\def\green{\color{green}}
\def\darkgreen{\color[rgb]{0,0.4,0.1}}
\def\darkblue{\color[rgb]{0,0.2,0.6}}
\def\skyblue{\color[rgb]{0.2,0.7,1.}}
%%control
\def\be{\begin{equation}}
\def\ee{\nonumber\end{equation}}
\def\bea{\begin{eqnarray*}}
\def\eea{\nonumber\end{eqnarray*}}
\def\bch{}
\def\ech{}
\def\bitem{\begin{itemize}}
\def\eitem{\end{itemize}}
\def\bcenter{\begin{center}}
\def\ecenter{\end{center}}
\def\bex{\begin{minipage}{0.2\textwidth}\includegraphics[width=0.6in]{jugelizi.png}\end{minipage}\begin{minipage}{0.76\textwidth}}
\def\eex{\end{minipage}}
\def\chtitle#1{\frametitle{\bch#1\ech}}
\def\bmat#1{\left(\begin{array}{#1}}
\def\emat{\end{array}\right)}
\def\bcase#1{\left\{\begin{array}{#1}}
\def\ecase{\end{array}\right.}
\def\bmini#1{\begin{minipage}{#1\textwidth}}
\def\emini{\end{minipage}}
\def\tbox#1{\begin{tcolorbox}#1\end{tcolorbox}}
\def\pfrac#1#2#3{\left(\frac{\partial #1}{\partial #2}\right)_{#3}}
%%symbols
\def\bropt{\,(\ \ \ )}
\def\sone{$\star$}
\def\stwo{$\star\star$}
\def\sthree{$\star\star\star$}
\def\sfour{$\star\star\star\star$}
\def\sfive{$\star\star\star\star\star$}
\def\rint{{\int_\leftrightarrow}}
\def\roint{{\oint_\leftrightarrow}}
\def\stdHf{{\textit{\r H}_f}}
\def\deltaH{{\Delta \textit{\r H}}}
\def\ii{{\dot{\imath}}}
\def\skipline{{\vskip0.1in}}
\def\skiplines{{\vskip0.2in}}
\def\lagr{{\mathcal{L}}}
\def\hamil{{\mathcal{H}}}
\def\vecv{{\mathbf{v}}}
\def\vecx{{\mathbf{x}}}
\def\vecy{{\mathbf{y}}}
\def\veck{{\mathbf{k}}}
\def\vecp{{\mathbf{p}}}
\def\vecn{{\mathbf{n}}}
\def\vecA{{\mathbf{A}}}
\def\vecP{{\mathbf{P}}}
\def\vecsigma{{\mathbf{\sigma}}}
\def\hatJn{{\hat{J_\vecn}}}
\def\hatJx{{\hat{J_x}}}
\def\hatJy{{\hat{J_y}}}
\def\hatJz{{\hat{J_z}}}
\def\hatj#1{\hat{J_{#1}}}
\def\hatphi{{\hat{\phi}}}
\def\hatq{{\hat{q}}}
\def\hatpi{{\hat{\pi}}}
\def\vel{\upsilon}
\def\Dint{{\mathcal{D}}}
\def\adag{{\hat{a}^\dagger}}
\def\bdag{{\hat{b}^\dagger}}
\def\cdag{{\hat{c}^\dagger}}
\def\ddag{{\hat{d}^\dagger}}
\def\hata{{\hat{a}}}
\def\hatb{{\hat{b}}}
\def\hatc{{\hat{c}}}
\def\hatd{{\hat{d}}}
\def\hatN{{\hat{N}}}
\def\hatH{{\hat{H}}}
\def\hatp{{\hat{p}}}
\def\Fup{{F^{\mu\nu}}}
\def\Fdown{{F_{\mu\nu}}}
\def\newl{\nonumber \\}
\def\vece{\mathrm{e}}
\def\calM{{\mathcal{M}}}
\def\calT{{\mathcal{T}}}
\def\calR{{\mathcal{R}}}
\def\barpsi{\bar{\psi}}
\def\baru{\bar{u}}
\def\barv{\bar{\upsilon}}
\def\qeq{\stackrel{?}{=}}
\def\torder#1{\mathcal{T}\left(#1\right)}
\def\rorder#1{\mathcal{R}\left(#1\right)}
\def\contr#1#2{\contraction{}{#1}{}{#2}#1#2}
\def\trof#1{\mathrm{Tr}\left(#1\right)}
\def\trace{\mathrm{Tr}}
\def\comm#1{\ \ \ \left(\mathrm{used}\ #1\right)}
\def\tcomm#1{\ \ \ (\text{#1})}
\def\slp{\slashed{p}}
\def\slk{\slashed{k}}
\def\calp{{\mathfrak{p}}}
\def\veccalp{\mathbf{\mathfrak{p}}}
\def\Tthree{T_{\tiny \textcircled{3}}}
\def\pthree{p_{\tiny \textcircled{3}}}
\def\dbar{{\,\mathchar'26\mkern-12mu d}}
\def\erf{\mathrm{erf}}
\def\const{\mathrm{constant}}
\def\pheat{\pfrac p{\ln T}V}
\def\vheat{\pfrac V{\ln T}p}
%%units
\def\fdeg{{^\circ \mathrm{F}}}
\def\cdeg{^\circ \mathrm{C}}
\def\atm{\,\mathrm{atm}}
\def\angstrom{\,\text{\AA}}
\def\SIL{\,\mathrm{L}}
\def\SIkm{\,\mathrm{km}}
\def\SIyr{\,\mathrm{yr}}
\def\SIGyr{\,\mathrm{Gyr}}
\def\SIV{\,\mathrm{V}}
\def\SImV{\,\mathrm{mV}}
\def\SIeV{\,\mathrm{eV}}
\def\SIkeV{\,\mathrm{keV}}
\def\SIMeV{\,\mathrm{MeV}}
\def\SIGeV{\,\mathrm{GeV}}
\def\SIcal{\,\mathrm{cal}}
\def\SIkcal{\,\mathrm{kcal}}
\def\SImol{\,\mathrm{mol}}
\def\SIN{\,\mathrm{N}}
\def\SIHz{\,\mathrm{Hz}}
\def\SIm{\,\mathrm{m}}
\def\SIcm{\,\mathrm{cm}}
\def\SIfm{\,\mathrm{fm}}
\def\SImm{\,\mathrm{mm}}
\def\SInm{\,\mathrm{nm}}
\def\SImum{\,\mathrm{\mu m}}
\def\SIJ{\,\mathrm{J}}
\def\SIW{\,\mathrm{W}}
\def\SIkJ{\,\mathrm{kJ}}
\def\SIs{\,\mathrm{s}}
\def\SIkg{\,\mathrm{kg}}
\def\SIg{\,\mathrm{g}}
\def\SIK{\,\mathrm{K}}
\def\SImmHg{\,\mathrm{mmHg}}
\def\SIPa{\,\mathrm{Pa}}
\def\secpage#1#2{\begin{frame}\bch\bcenter{\bf \Huge #1} \skipline \tbox{\bcenter #2\ecenter}\ecenter\ech\end{frame}}

\graphicspath{{figure/}}
\usepackage{enumerate}
%\usepackage{fancyhdr}
%\cpic{<尺寸>}{<文件名>}}用于生成居中的图片。
\newcommand{\cpic}[2]{
\begin{center}
\includegraphics[scale=#1]{#2}
\end{center}
}
%\cpicn{<尺寸>}{<文件名>}{<注释>}用于生成居中且带有注释的图片，其label为图片名。
\newcommand{\cpicn}[3]
{
\begin{figure}[H]
\cpic{#1}{#2}
\caption{\color{red}#3\label{#2}}
\end{figure}
}

\newtheorem{definition}{\hspace{2em} 定义}[section]
\newtheorem{conclusion}{\hspace{2em} 重要结论}[section]
\newtheorem{practice}{\hspace{2em} 小练习}[section]
\title{$SU(2)$群的奇异特性}
\author{徐昊霆}
\begin{document}
\maketitle
\tableofcontents
\section{$SU(N)$群复习}
\subsection{定义、张量方法构造表示}
\begin{definition}[$SU(N)$ 群]
  同时满足$U^{\dagger}U = 1$且$\mathrm{det} U = 1$的群。
\end{definition}
回顾$SU(N)$群对矢量和张量的作用，对于矢量
\beq
\psi^i \rightarrow u^{ij}\psi^{j}
\eeq
对于具有$m$个指标的张量
\beq
\psi^{i_1\cdots i_m} \rightarrow U^{i_1j_1}\cdots U^{i_mj_m} \psi^{j_1\cdots j_m}
\eeq
一般来说只需要考虑对称张量和全反对称张量。于是，对于$SU(2)$群，张量$\psi^{i_1\cdots i_m}$有$m+1$个独立分量，因此对应的表示是$m+1$维表示。之后我们计算任何一个张量的迹，我们发现它不是像$SO(N)$群一样像一个张量一样独立变换。为了让迹独立变换，我们定义了上下指标，上下指标的变换为
\bea
\psi^{i} &\rightarrow & U^i_j \psi^j\\
\psi_{i} &\rightarrow & \left(U^{\dagger}\right)^i_j \psi^j\\
\eea
于是可以知道一个张量究竟是怎样变换的。再利用$\mathrm{det}U =1$，我们得到
\beq\label{invariant}
\mathrm{det} U = \epsilon_{i_1i_2\cdots i_N} U^{i_1}_{j_1}\cdots U^{i_1}_{j_1} =\epsilon_{j_1j_2\cdots j_N}
\eeq
因此，在变换下，可以得到下面的等式\footnote{这里取一个特殊的例子，有$m$个指标的写起来过于繁杂}
\beq
\phi_{ik} = \epsilon_{ipq} \phi^{pq}_k
\eeq
利用式~\ref{invariant}可以验证，在$U$的变换下上面的等式确实成立。因此得到下面的重要结论
\begin{conclusion}
  在$SU(N)$群中，我们使用$\epsilon$符号升降张量的指标。
\end{conclusion}
\subsection{$SU(N)$代数}
考虑无穷小变化
\beq
H = I + i\epsilon H
\eeq
并利用关系$U^{\dagger}U = I$，可以得到$H$是{\color{red} 无迹厄米矩阵}。$N$维无迹厄米矩阵有$(N-1)+N(N-1) = N^2-1$个分量。对于$SU(2)$群，无迹厄米矩阵的集合就是三个泡利矩阵。我们回忆$SO(N)$群的生成元有$\frac{1}{2}N(N-1)$个，于是我们惊奇的发现$2^2 -1 = \frac{1}{2} 3\times 2 =3$，暗示了$SU(2)$群和$SO(3)$群惊人的联系。
\section{$SU(2)$和$SO(3)$局域同构}
在这一节我们将看到，一个$SU(2)$中有两个$SO(3)$。我们想办法把三维空间的转动对应到$SU(2)$中的元素上去。我们把三维空间的矢量$\vec{x}$先映射到$SU(2)$的无迹厄米矩阵$X$上去，并令
\beq
X = x\sigma_1 + y\sigma_2 +z\sigma_3 =
\begin{pmatrix}
  z    & x-iy \\
  x+iy & z  
\end{pmatrix}
\eeq
一会我们就将看到，如果这样映射，就可以把两个群对应起来。我们注意到$X$的行列式为
\beq
\mathrm{det} X = -(x^2+y^2+z^2) = -\vec{x}^2
\eeq
现在我们做变换$X^{\prime} = U^{\dagger}XU$，我们会发现变换后的$X^{\prime}$仍然是一个无迹厄米矩阵，且它的行列式不变。先求迹
\beq
\mathrm{tr} X^\prime = \mathrm{tr} XUU^{\dagger} = \mathrm{tr}X = 0
\eeq
上面的推导需要用一点张量指标的小技巧。也很容易证明它的厄米性。现在，新的$X^{\prime}$又可以用一组新的系数来代表，它们是$(x^\prime,y^{\prime},z^\prime)$，我们计算矩阵$X^{\prime}$的行列式，我们会得到$\mathrm{det}X = \mathrm{det}X^\prime$，我们发现，这种变换也是保证模长不变的线性变换,根据定义，这正是一个三维空间的旋转。另外，映射$f:U\rightarrow R$也构成群。

值得注意的是，映射$f:U\rightarrow R$是一个2对1的映射，因为$-U$也可以映射到同一个$R$上，因为$(-U)^{\dagger} X (-U)= U^\dagger XU$，在$SU(2)$群中，$U$和$-U$显然不可能是一个东西。所以，我们又有了重要结论
\begin{conclusion}
  $SU(2)$ covers $SO(3)$ twice.
\end{conclusion}

上面我们只是大概讨论了一下做那样的变化你可以把$SU(2)$映射到$SO(3)$,下面我们来直接计算一下。我们知道$SU(2)$的元素可以表示为
\beq
U = e^{i\varphi_a \sigma_a/2}
\eeq
我们定义$\vec{\varphi} = \varphi \hat{\varphi}$，利用
\beq
\left(\vec{\varphi}\cdot \vec{\sigma}\right)^2 = \varphi^2
\eeq
于是直接将群元直接展开，并将奇偶项拆开
\beq
\begin{aligned}
  U &=& \sum_{n=0}^{\infty} \frac{i^n}{n!}\left(\frac{\vec{\varphi}\cdot \vec{\sigma}}{2}\right)^n \\
  &=& \cos \frac{\varphi}{2} + i \hat{\varphi}\cdot \vec{\sigma} \sin\frac{\varphi}{2}
\end{aligned}
\eeq
现在直接计算$U^\dagger X U$，并利用泡利矩阵的性质$\sigma_a\sigma_b = \delta_{ab}I+ i\epsilon_{abc}\sigma_c$，如果对于$z$轴的旋转，就会得到原来我们熟悉的变换。

根据刚刚的计算，如果是绕$z$轴旋转，那么有$U(\varphi)=e^{i\varphi \sigma_3/2}$，如果我们考虑$2\pi$的旋转，我们得到
\beq
U(2\pi) = -I
\eeq
这里又反映了我们之前的重要结论，$SU(2)$中包含两个$SO(3)$，在下一节我们会对这个问题有更深刻的理解。
\section{$SU(2)$的李代数}
在之前我们讲过，$SU(2)$的李代数与$SO(3)$有同样的结构，他们满足对易关系
\beq
\left[T^i,T^j\right] = i\epsilon^{ijk} T^k
\eeq
于是同样的可以定义升降算符，满足$T^\pm = T^1\pm T^2$，于是对易关系变为
\beq
\left[T^3, T^\pm\right] = \pm T^{\pm},\spa \left[ T^+,T^-\right] = 2T^3
\eeq
于是选择$T^3$的本征态，我们又可以重复之前李代数的操作。我们将看到$SU(2)$的表示是$2j+1$维的，其中$j$既可以取整数又可以取半整数。我们其实已经可以看出为什么了，这是因为具有$m$个指标的张量贡献一个$m+1$维表示，直接令$m=2j$，就可以得到上面的半整数了。但是，这样说到目前为止还有一点不严谨，因为我们知道，对于一般$SU(N)$群，上标的张量和下标的张量是不一样的，那么下标会不会贡献其他的表示?
\section{$SU(2)$的奇妙性质}
\subsection{只需考虑上标对称张量}
我们之前学习$SO(3)$的时候，我们知道了$SO(3)$的特殊性，即我们只需要考虑$n$维无迹对称张量就行了，其他的$SO(N)$群则没有这样的特性。

根据刚刚的指标升降公式，我们可以把下标的全部都升上去，例如对张量$T^{pqijk} = \epsilon ^{pm}\epsilon^{qn} T_{mn}^{ijk}$。我们也只需要考虑对称的张量，因为如果有一个张量，比如说，4个指标的张量$T^{ijkl}$，总可以对两个指标$i,k$构造对称部分和反对称部分
\bea
S^{ijkl} &=& T^{ijkl} + T^{kjil} \\
A^{ijkl} &=& T^{ijkl} - T^{kjil}
\eea
然后对于反对称张量，可以构造一个二阶的对称张量$\epsilon_{ik} A^{ijkl}$，所以高维的反对称张量总可以变成低维的张量，最后全变成一维的张量，而一维的张量没有反对称和对称这一说法。所以说，为了构造$SU(2)$群的不可约表示，我们只需要考虑上标的对称张量即可。这和其他的群是不同的。因此，我们就可以说对于$m$个指标的张量，我们只需要考虑对称张量就完了，所以得到$SU(2)$有$m+1$维的表示。
\begin{conclusion}
  $SU(2)$有$m+1$维的表示。
\end{conclusion}
\begin{practice}
  试证明上面的论述对于$SU(3)$群不成立。
\end{practice}
所以，我们看到了之前李代数的半整数的来源了，是$SU(2)$群。

\subsection{$SU(2)$ is pseudoreal}
在上一讲我们提到，必须严格区分上下标，但是我们刚刚又讲过，不需要下标的东西。这不是矛盾吗?不需要下标到底是啥意思，它们不是差一个共轭吗?我们回想到我们在第二章学过一个群表示的real 和 pseudoreal，实际上，$SU(2)$是pseudoreal的。回顾一下定义
\begin{definition}
  一个群的表示是pseudoreal的，如果存在一个矩阵$S$，使得
  \beq
  D(g)^* = SDS^{-1}
  \eeq
\end{definition}
可以验证，这个矩阵$S$可以是$\sigma_2$。
\begin{practice}
  验证矩阵$S$可以是$\sigma_2$
\end{practice}
\begin{conclusion}
  $SU(2)$ is pseudoreal.
\end{conclusion}

\section{$SU(N)$和$U(N)$的关系}
$U(N)$群有两种，第一种是$U(1)$群，它的形式为$e^{i\varphi}I$，第二种就是我们所说的$SU(N)$群，聪明的物理学家这时候路过，写下$U(N) = SU(N)\otimes U(1)$。

但是这是错的。

我们没有意识到，这两种是有交集的。考虑下面这些元素
\be
e^{i2\pi k/N} I, \,\spa k=1,\cdots,N
\ee
你惊奇的发现，这个矩阵的行列式为1，而且它还是$U(1)$群的元素。这些群是啥?$Z_N$群的元素！所以正确的写法为
\begin{conclusion}
\beq
U(N) = \left(SU(N)/Z_N\right)\otimes U(1)
\eeq
\end{conclusion}
上面的写法如果对于$N=2$，又一次暗示了$SU(2)$包含了两个$SO(3)$。我们回想起强力、弱力和电磁相互作用的对称群是$SU(3)\oplus SU(2)\oplus U(1)$，这个写法应该也不严谨。但是在物理中，整体的对称性没有出现，而只是李代数或者局域不变性出现在拉格朗日量中。

%\bibliographystyle{siam}
%\bibliography{cites}
\end{document}
